% Part: applied-modal-logics
% Chapter: temporal-logic
% Section: properties-accessibility

\documentclass[../../../include/open-logic-section]{subfiles}

\begin{document}

\olfileid{aml}{tl}{acc}

\olsection{سمے دے فریماں دیاں خاصیتاں}

ساڈے سمے دے نمونے تعلق~$\prec$ اُتے کوئی شرط نہیں لاؤندے، ایس لئی
ساڈے جانے پہچانے مسلّمیاں وچوں صرف~$K$، یا اوہدے برابر دے
$K_{\Gtemp}$ تے~$K_{\Htemp}$، سارے نمونیاں وچ لاگو ہوندے نیں:
\begin{align*}
\tag{$K_{\Gtemp}$}  \Gtemp (p \to q) & \to (\Gtemp p \to \Gtemp q)\\
\tag{$K_{\Htemp}$}  \Htemp (p \to q) & \to (\Htemp p \to \Htemp q)
\end{align*}

پر جے اسیں چاہندے آں کہ ساڈے نمونے سماں وکھاؤن اُتے ہور سخت شرطاں لاؤن،
مثلاً ایہ یقینی بنائیے کہ $\prec$ اک سیدھی ترتیب اے، تاں ساڈے
نمونیاں وچ ہور بیان وی معتبر ہون گے۔

\begin{table}[t]
  \begin{tabular}{| p{.48\textwidth} || p{.48\textwidth} |}
    \hline
    {\emph{جے $\prec$ \dots اے}} & {\emph{تاں \dots،~$\mModel{M}$ وچ سچ اے}} \\
    \hline \hline
    \emph{متعدی}: \newline
    $\forall u \forall v \forall w ((u \prec v \land v \prec w) \lif u \prec w)$ & 
    $\Ftemp \Ftemp p \lif \Ftemp p$  \\
    \hline 
    \emph{سیدھی ترتیب}: \newline
    $\forall w \forall v (w \prec v \lor w = v \lor v \prec w)$ &  
    $(\Ftemp \Ptemp  p \lor \Ptemp \Ftemp p) \lif (\Ptemp p \lor p \lor \Ftemp p)$\\
    \hline
    \emph{گنجان}: \newline
    $\forall w \forall v (w \prec v \to \exists u(w \prec u \land u \prec v))$ &  
    $\Ftemp p \lif \Ftemp \Ftemp p$ \\
    \hline
    \emph{لنگھے پاسے بے حد}: \newline
    $\forall w \exists v( v \prec w)$ &  
    $\Htemp p \to \Ptemp p$ \\
    \hline
    \emph{آؤن والے پاسے بے حد}: \newline
    $\forall w \exists v( w \prec v)$ &  
    $\Gtemp p \to \Ftemp p$ \\
    \hline
  \end{tabular}
  \caption{سمے دے فریماں دیاں کجھ خاصیتاں تے اوہناں دے برابر فارمولے۔}
  \ollabel{tab:correspondence}
\end{table} 

\olref{tab:correspondence} دیاں کئی خاصیتاں سمے نوں وکھاؤن والے
نمونے لئی چنگیاں لگ سکدیاں نیں۔ پر ایہ گل دھیان وچ رکھو کہ بھاویں
اسیں تعلق $\prec$ اُتے من چاہیاں شرطاں لا سکدے آں، ساریاں شرطاں دے
برابر ایسے !!{formula}s نہیں ہوندے جیہڑے سمے دی منطق دی زبان وچ لکھے
جا سکن۔ مثلاً غیر انعکاسی ہون دی شرط، یعنی
$\forall w \lnot (w \prec w)$، دے برابر سمے دی منطق دا کوئی فارمولا نہیں۔

\end{document}
